\documentclass[11pt]{article}

\usepackage[margin=1in]{geometry}
\usepackage{booktabs}
\usepackage{enumitem}
\usepackage{hyperref}
\usepackage{xcolor}

\title{Curling Experiment Design}
\author{}
\date{}

\begin{document}
\maketitle

\section{Overview}

This document describes the benchmark plan for evaluating Curling's
control-plane/data-plane separation.  Curling decouples request admission from
private-side execution by placing a queue between the API service and the
private executor.  The executor can either process requests one by one or apply
batch scheduling with artifact reuse.  The experiments compare this decoupled
design against coupled baselines in which the API waits for private-side work to
finish synchronously.

The benchmark uses fixed traces generated from workload configuration files.
Each execution mode receives the same request sequence, object size mix, and
request priority distribution.  This makes the comparison focus on execution
mode behavior rather than trace randomness.

\section{Main Questions}

The experiments are designed to answer the following questions.

\begin{enumerate}[leftmargin=*]
  \item \textbf{What is the overhead of decoupling?}
  When reuse is disabled, how much latency or throughput cost does the queue
  introduce compared with coupled synchronous baselines?

  \item \textbf{When does reuse pay off?}
  Under read-heavy workloads with repeated targets, how much does decoupled
  batch scheduling and artifact reuse improve completion latency, throughput,
  and private-side execution work?

  \item \textbf{How robust is the design across workload mixes?}
  Does Curling remain competitive for mixed and write-heavy workloads where
  reuse opportunities are weaker?

  \item \textbf{How does load affect each mode?}
  How do admission latency, completion latency, and throughput change as client
  concurrency increases?

  \item \textbf{How does object size affect performance?}
  Do large objects amplify the benefit of reuse or the cost of synchronous data
  movement through the API service?
\end{enumerate}

\section{Systems Under Test}

The same deployment is configured in one mode at a time.  Object storage is
used in all modes; the key difference is whether the API waits synchronously for
private execution or enqueues the request for asynchronous execution.

\begin{table}[h]
\centering
\caption{Execution modes and baselines.}
\label{tab:modes}
\begin{tabular}{llll}
\toprule
\textbf{Mode} & \textbf{API behavior} & \textbf{Private behavior} & \textbf{Reuse} \\
\midrule
Coupled RPC baseline &
Synchronous RPC call &
Executor runs request immediately &
No \\
Coupled storage API baseline &
Synchronous storage API access &
Storage service reads/writes files &
No \\
Decoupled baseline &
Enqueue request &
Queue worker runs requests one by one &
No \\
Decoupled reuse &
Enqueue request &
Batch scheduler with artifact reuse &
Yes \\
\bottomrule
\end{tabular}
\end{table}

\subsection{Baseline Rationale}

\textbf{Coupled RPC baseline.}
This baseline represents a tightly coupled design where the API service calls a
private executor directly and waits for the request to finish.  It avoids queue
overhead but exposes the API to private-side execution latency.

\textbf{Coupled storage API baseline.}
This baseline represents a coupled design where the private node exposes a
storage API and the control-plane API mediates data movement.  It is useful for
measuring the cost of moving data through the API service.

\textbf{Decoupled baseline.}
This mode keeps the queue-based architecture but disables reuse.  It isolates
the cost and benefit of asynchronous decoupling from the benefit of reuse.

\textbf{Decoupled reuse.}
This is the full Curling design.  The executor polls pending requests in
batches, first binds requests to non-expired artifacts when possible, then
groups remaining read requests by reuse key.  The current reuse key is the read
target.  Write requests are scheduled as single-request execution units.  Units
are ordered by scheduling score; the current implementation primarily uses
request priority.

\section{Workloads}

Workloads are defined as JSON files under \texttt{experiments/workloads/}.
The dataset generator reads a workload file and emits a trace JSON file plus
metadata.  The private node materializes read targets from the trace, while the
client node materializes write input files from the same trace.

\begin{table}[h]
\centering
\caption{Default workload presets.}
\label{tab:workloads}
\begin{tabular}{lllll}
\toprule
\textbf{Workload} & \textbf{Read/write mix} & \textbf{Read targets} &
\textbf{Repeat mode} & \textbf{Purpose} \\
\midrule
\texttt{read-heavy-reuse} &
90\% read, 10\% write &
10 &
Zipf, $\alpha=1.2$ &
Highlight artifact reuse \\
\texttt{mixed} &
50\% read, 50\% write &
100 &
Zipf, $\alpha=0.8$ &
Model normal exchange traffic \\
\texttt{write-heavy-no-reuse} &
10\% read, 90\% write &
100 &
Unique &
Measure decoupling when reuse is weak \\
\bottomrule
\end{tabular}
\end{table}

\subsection{Object Size Distribution}

Each request is assigned an object size label.  The default profile is shown in
Table~\ref{tab:sizes}.  Unless otherwise stated, experiments use the default
mix: 70\% small, 20\% medium, and 10\% large.

\begin{table}[h]
\centering
\caption{Default object size profile.}
\label{tab:sizes}
\begin{tabular}{lll}
\toprule
\textbf{Label} & \textbf{Bytes} & \textbf{Default ratio} \\
\midrule
Small & 1,024 & 0.7 \\
Medium & 1,048,576 & 0.2 \\
Large & 10,485,760 & 0.1 \\
\bottomrule
\end{tabular}
\end{table}

\subsection{Priority Policies}

The trace records a priority for every request.  The default policy is
\texttt{uniform}, where all requests have priority 0.  A second policy,
\texttt{small-first}, can assign higher priority to smaller objects, for
example small:2, medium:1, large:0.  Priority experiments should use the same
trace across all modes so that only scheduling behavior differs.

\section{Metrics}

The benchmark records per-request and aggregate metrics.

\begin{itemize}[leftmargin=*]
  \item \textbf{Admission latency}: time spent by the submit call.
  \item \textbf{Completion latency}: time from submit start to final state.
  \item \textbf{Operation latency}: full client-side operation latency,
  including staged upload for writes.
  \item \textbf{Throughput}: successfully completed requests per second.
  \item \textbf{Final state distribution}: number of requests ending in states
  such as released, committed, failed, or timed out.
  \item \textbf{Mode-specific behavior}: artifact reuse hits, queue delay, and
  execution-unit batching should be collected from logs or future counters when
  available.
\end{itemize}

Latency should be reported with at least median, p95, and p99.  Results should
also separate read and write requests because the two paths exercise different
parts of the system.

\section{Experimental Parameters}

The main sweep parameters are listed in Table~\ref{tab:params}.  A compact
default matrix is sufficient for early runs; larger runs should increase request
count and repeat each point multiple times.

\begin{table}[h]
\centering
\caption{Benchmark parameters.}
\label{tab:params}
\begin{tabular}{lll}
\toprule
\textbf{Parameter} & \textbf{Default values} & \textbf{Reason} \\
\midrule
Execution mode &
Four modes in Table~\ref{tab:modes} &
Compare baselines and full design \\
Workload &
Three presets in Table~\ref{tab:workloads} &
Cover reuse-heavy, mixed, and low-reuse cases \\
Concurrency &
1, 10, 50 &
Measure behavior from light to high load \\
Requests per run &
1,000 default, larger for final runs &
Balance stability and runtime \\
Object size mix &
70/20/10 small/medium/large &
Represent heterogeneous object sizes \\
Priority policy &
Uniform, optional small-first &
Evaluate scheduling sensitivity \\
Random seed &
Fixed per trace &
Make modes comparable \\
\bottomrule
\end{tabular}
\end{table}

\section{Experiment Matrix}

\subsection{Experiment 1: Decoupling Overhead}

\textbf{Purpose.}
Measure the overhead introduced by queue-based decoupling when reuse is not the
dominant factor.

\textbf{Workload.}
\texttt{write-heavy-no-reuse}.

\textbf{Modes.}
Coupled RPC baseline, coupled storage API baseline, and decoupled baseline.

\textbf{Procedure.}
Generate one fixed trace from \texttt{write\_heavy\_no\_reuse.json}.  Materialize
private read files on the private node and write input files on the client node.
Run the same trace for each mode at concurrency 1, 10, and 50.

\textbf{Expected interpretation.}
The gap between coupled RPC and decoupled baseline estimates queue and
asynchronous execution overhead.  The gap between coupled RPC and coupled
storage API estimates the cost of API-mediated data transfer.

\subsection{Experiment 2: Reuse Benefit}

\textbf{Purpose.}
Quantify the benefit of artifact-level and intra-batch reuse.

\textbf{Workload.}
\texttt{read-heavy-reuse}.

\textbf{Modes.}
Decoupled baseline and decoupled reuse.  Coupled baselines may be included as
additional references.

\textbf{Procedure.}
Generate one fixed read-heavy trace.  Run it against decoupled baseline and
decoupled reuse at concurrency 1, 10, and 50.  Repeat with a larger request
count if p95 and p99 latency are unstable.

\textbf{Expected interpretation.}
The full design should reduce private-side work and improve completion latency
when many read requests share a small target set.  Admission latency may remain
similar between the two decoupled modes because both use the API queue path.

\subsection{Experiment 3: Mixed Workload Robustness}

\textbf{Purpose.}
Evaluate whether reuse scheduling helps under a more balanced read/write mix
without harming write-heavy behavior.

\textbf{Workload.}
\texttt{mixed}.

\textbf{Modes.}
All four modes.

\textbf{Procedure.}
Generate one fixed mixed trace and run it at each concurrency level.  Report
read and write latencies separately as well as aggregate throughput.

\textbf{Expected interpretation.}
The decoupled reuse mode may still benefit repeated reads, but writes should not
be starved.  If write p95 or p99 latency grows substantially, the scheduling
score or tenant budget policy needs adjustment.

\subsection{Experiment 4: Object Size Sensitivity}

\textbf{Purpose.}
Determine whether benefits and overheads scale with object size.

\textbf{Workload.}
Use \texttt{read-heavy-reuse} and \texttt{write-heavy-no-reuse}; vary the object
size mix.

\textbf{Parameter sweep.}
Run three mixes: small-only, balanced, and large-heavy.  For example:
\texttt{small:1.0,medium:0,large:0},
\texttt{small:0.7,medium:0.2,large:0.1}, and
\texttt{small:0.2,medium:0.3,large:0.5}.

\textbf{Expected interpretation.}
Large objects should increase the cost of repeated data processing and data
movement, making reuse more valuable in read-heavy workloads and making
API-mediated transfer more expensive in coupled storage API mode.

\subsection{Experiment 5: Priority Scheduling}

\textbf{Purpose.}
Check whether the scheduler can favor high-priority requests without severely
hurting low-priority throughput.

\textbf{Workload.}
Use \texttt{mixed} or \texttt{read-heavy-reuse} with
\texttt{priority\_policy=small-first}.

\textbf{Modes.}
Decoupled baseline and decoupled reuse.

\textbf{Procedure.}
Generate two traces from the same workload: one with uniform priority and one
with small-first priority.  Compare latency by priority class and object size.

\textbf{Expected interpretation.}
In decoupled reuse mode, high-priority execution units should see lower latency.
Low-priority requests should still make progress; otherwise tenant budgets or
aging terms should be strengthened.

\section{Execution Procedure}

For each workload and parameter point, use the following procedure.

\begin{enumerate}[leftmargin=*]
  \item Generate a trace on the admin machine:
\begin{verbatim}
python experiments/benchmark/generate_dataset.py \
  --role admin \
  --workload-json experiments/workloads/read_heavy_reuse.json \
  --output-dir experiments/benchmark/generated/read-heavy-reuse \
  --seed 1
\end{verbatim}

  \item Materialize private read files on the private node:
\begin{verbatim}
python experiments/benchmark/generate_dataset.py \
  --role private \
  --trace-json experiments/benchmark/generated/read-heavy-reuse/read-heavy-reuse.trace.json \
  --private-export-dir private-data/exports
\end{verbatim}

  \item Materialize client write input files on the client node:
\begin{verbatim}
python experiments/benchmark/generate_dataset.py \
  --role client \
  --trace-json experiments/benchmark/generated/read-heavy-reuse/read-heavy-reuse.trace.json \
  --write-input-dir experiments/benchmark/generated/read-heavy-reuse/writes
\end{verbatim}

  \item Configure one execution mode in \texttt{inventory.env}, regenerate
  environment files, and restart node1 and node3.

  \item Run the benchmark with the generated trace:
\begin{verbatim}
bash deploy/scripts/run_benchmark.sh decoupled-reuse \
  --trace-json experiments/benchmark/generated/read-heavy-reuse/read-heavy-reuse.trace.json \
  --trace-base-dir experiments/benchmark/generated/read-heavy-reuse \
  --concurrency 10
\end{verbatim}

  \item Save the CSV and summary JSON from
  \texttt{experiments/benchmark/results/}.  Record the git commit, environment
  configuration, node types, network placement, and start time.
\end{enumerate}

\section{Reporting}

Each experiment should report:

\begin{itemize}[leftmargin=*]
  \item throughput by mode and concurrency;
  \item admission, completion, and operation latency by mode;
  \item read-only and write-only latency breakdowns;
  \item p50, p95, and p99 latency;
  \item final state counts and failure details;
  \item for reuse experiments, artifact hit rate and average execution-unit
  size when counters are available.
\end{itemize}

Plots should use the same y-axis scale when comparing modes within one
experiment.  For latency plots, use log scale when tail latency spans more than
one order of magnitude.

\section{Validity Considerations}

The following factors should be controlled or reported.

\begin{itemize}[leftmargin=*]
  \item Use the same trace for all modes in a comparison.
  \item Warm up services before measuring, or report cold-start results
  separately.
  \item Clear or isolate artifacts between runs when measuring first-run reuse
  behavior.
  \item Keep object storage, database, and queue services configured identically
  across modes.
  \item Repeat final measurements at least three times and report mean and
  variance.
  \item Separate API admission latency from end-to-end completion latency,
  because decoupled modes intentionally return earlier on admission.
\end{itemize}

\end{document}
